% Hexagon note, 23 Nov 2026. CC0 1.0.
\documentclass[11pt,a4paper]{article}
\usepackage[T1]{fontenc}
\usepackage{lmodern}
\usepackage{microtype}
\usepackage[margin=2.7cm]{geometry}
\usepackage{amsmath,amssymb,amsthm}
\usepackage{booktabs}
\PassOptionsToPackage{hyphens}{url}
\usepackage[hidelinks]{hyperref}
\emergencystretch=2em

\theoremstyle{plain}
\newtheorem{theorem}{Theorem}
\newtheorem{lemma}[theorem]{Lemma}
\newtheorem{proposition}[theorem]{Proposition}
\newtheorem{corollary}[theorem]{Corollary}
\theoremstyle{remark}
\newtheorem{remark}[theorem]{Remark}
\newtheorem{example}[theorem]{Example}

\newcommand{\Q}{\mathbf Q}
\newcommand{\Z}{\mathbf Z}
\newcommand{\C}{\mathbf C}
\newcommand{\Ag}{\mathcal A_g}
\newcommand{\OO}{\mathcal O}
\newcommand{\Gal}{\operatorname{Gal}}
\newcommand{\End}{\operatorname{End}}
\newcommand{\disc}{\operatorname{disc}}
\newcommand{\Qb}{\overline{\mathbf Q}}
\newcommand{\workbenchurl}{https://github.com/KokunoYumeto/mathematics-workbenches}
\newcommand{\recordurl}{https://github.com/KokunoYumeto/mathematics-workbenches}

\title{Galois orbits of CM points in $\mathcal A_g$ with an explicit exponent,\\ after Tsimerman and Gaudron--R\'emond}
\author{Entered at deposit}
\date{23 November 2026 \quad CC0 1.0}

\begin{document}
\maketitle

\begin{abstract}
\noindent
Tsimerman proved the Andr\'e--Oort conjecture for the moduli space $\Ag$ of principally polarized abelian varieties
by showing that the Galois orbit of a CM point is at least a power of the discriminant of its endomorphism algebra;
the exponent was not specified and the constant was not effective. We prove that every CM point $x$ of $\Ag$ satisfies
\[
 |\Gal(\Qb/\Q)\cdot x|\ \geq\ c_g\,|\disc R_x|^{1/(4g)},
\]
where $R_x$ is the centre of the endomorphism ring of $A_x$ and $c_g>0$ is an explicit constant; a sharper logarithmic
form has exponent $1/(2g)$. The bound holds for every centre order, maximal or not, and for products. The proof
combines the 2023 isogeny theorems of Gaudron and R\'emond, which bound the discriminant of the $\ell$-adic Galois
images, with the averaged Colmez formula for heights. Alp\"oge's theorem on Dirichlet $L$-functions is used once, to
exclude an exceptional real zero in the height estimate; the recent unconditional zero-exclusion results of Banks
(after OpenAI), Cureton--Thorner and Liu would serve as well, with other constants. A rational CM point of
$\mathcal A_2$ has centre discriminant less than $10^{80}$. The mathematics and proofs were produced by GPT-6 Astra
(ChatGPT Pro mode); they were checked step by step by Claude Opus 5.5 at effort high, which also wrote this text. No
human has checked them.
\end{abstract}

\section{The theorem}

Let $\Ag/\Q$ be the coarse moduli space of principally polarized abelian varieties of dimension $g$. A point
$x\in\Ag(\Qb)$ is a \emph{CM point} if the abelian variety $A_x$ it represents has complex multiplication, that is,
if $\End^0(A_x)$ contains a commutative semisimple $\Q$-algebra of dimension $2g$. Put
\[
 R_x=Z(\End_{\Qb}A_x),\qquad D_x=|\disc R_x|,
\]
the discriminant of the centre with respect to its own trace form. The algebra $R_x\otimes\Q$ is a product of CM
fields, one for each isogeny class of simple factors of $A_x$, and $R_x$ is an order in it, not necessarily maximal.
Thus $D_x\geq3$. The degree $[\Q(x):\Q]$ of the field generated by the coordinates of $x$ is the cardinality
$|\Gal(\Qb/\Q)\cdot x|$ of its Galois orbit.

Throughout, $\mathrm e=\exp(1)$, logarithms are natural, and for positive integers $r$ and $g$ we put
\begin{align*}
 a_r&=2^{r-1}\Bigl(\tfrac{75}2+16\,r!\Bigr),& b_r&=4r(2^r-1),&
 A_g&=\max_{r\leq g}a_r,& B_g&=\max_{r\leq g}b_r,\\
 C_g&=241(\mathrm eg)^{2g}g^5,& s_g&=2\cdot3^{4g^2},
\end{align*}
\[
 W_g=g\bigl(A_g+\tfrac12\bigr)+\frac{gB_g+\tfrac32g+1}{\log3},\qquad
 c_g=\frac{\mathrm e}{4g\,s_gC_gW_g}.
\]

\begin{theorem}\label{thm:main}
Every CM point $x\in\Ag(\Qb)$ satisfies
\[
 |\Gal(\Qb/\Q)\cdot x|\ \geq\ \frac{D_x^{1/(2g)}}{s_gC_gW_g\log D_x}\ \geq\ c_g\,D_x^{1/(4g)}.
 \tag{1.1}
\]
More precisely, if $A$ is an abelian variety of dimension $g$ with complex multiplication over a number field $K$
over which all its geometric endomorphisms are defined, and $D=|\disc Z(\End A)|$, then
\[
 [K:\Q]\ \geq\ \frac{D^{1/(2g)}}{C_gW_g\log D}.
 \tag{1.2}
\]
\end{theorem}

\begin{table}[ht]
\centering
\begin{tabular}{@{}cccc@{}}
\toprule
$g$ & exponent $1/(4g)$ & $\log_{10}c_g$ & $W_g$\\
\midrule
1&$1/4$&$-7.405$&$59.92$\\
2&$1/8$&$-17.746$&$326.33$\\
3&$1/12$&$-31.623$&$1837.89$\\
4&$1/16$&$-49.447$&$14370.20$\\
\bottomrule
\end{tabular}
\caption{The constants for small $g$. The proofs use only the exact formulas.}
\end{table}

\paragraph{What the theorem says.}
A CM point of $\Ag$ whose endomorphisms have a centre of large discriminant cannot be defined over a field of small
degree: a point of degree at most $d$ has $D_x\leq(d/c_g)^{4g}$. In particular, for each $g$ and $d$, the CM points
of degree at most $d$ form a finite set whose centre discriminants are bounded explicitly. For example:

\begin{corollary}\label{cor:surfaces}
Every CM point of $\mathcal A_2$ defined over $\Q$ has $D_x<10^{80}$.
\end{corollary}

For $g=1$ the theorem is weak: there $D_x$ is the discriminant $|\Delta|$ of the order of $A_x$, and the orbit has
exactly $h(\Delta)$ elements, a class number. For $|\Delta|=163$ and $427$ the orbits have $1$ and $2$ elements,
while the middle term of (1.1) is about $1.45\cdot10^{-7}$ and $1.97\cdot10^{-7}$. The uniform constants are crude.
The point of the theorem is that they are explicit, in every dimension.

\section{Why it matters, and whose theorems carry it}

\paragraph{The question.}
Pila and Zannier's strategy reduces the Andr\'e--Oort conjecture for $\Ag$ to two inputs. The first is a count of
rational points on definable sets; the second is a lower bound for Galois orbits of CM points, polynomial in the
discriminant. Edixhoven had suggested the second as a conjecture. Tsimerman proved it \cite[Theorem~5.1]{Ts18}: for
some $b_g>0$, $D_x\ll_g|\Gal(\Qb/\Q)\cdot x|^{b_g}$. With the counting theorem this gave the Andr\'e--Oort conjecture
for $\Ag$ \cite[Theorem~5.2]{Ts18}. In Tsimerman's proof the exponent is any number below $1/(4c_g)$, where $c_g$ is the
exponent of the Masser--W\"ustholz isogeny theorem \cite[proof of Theorem~4.1]{Ts18}, and the implied constant is not
effective: the relative class-number bound $|\mathrm{Cl}(E)|/|\mathrm{Cl}(E^+)|\gg D_E^{1/4-o(1)}$ and the height bound
both come from the Brauer--Siegel theorem \cite[\S\S1.1, 3]{Ts18}. The ineffectivity is Siegel's: a possible real
zero of a quadratic Dirichlet $L$-function very close to $1$.

\paragraph{What is proved here.}
Theorem~\ref{thm:main} gives the orbit bound for all CM points of $\Ag$, not only those with maximal endomorphism
order, with the explicit exponent $1/(4g)$ and an explicit constant. The exponent comes from the isogeny theorems of
Gaudron and R\'emond \cite{GR23}, which bound the discriminant of the $\ell$-adic Galois images of an abelian variety
by a power $2n\leq 2g$ of its degree times its height. The height comes from the averaged Colmez formula, proved by
Yuan and Zhang \cite{YZ} and by Andreatta, Goren, Howard and Madapusi Pera \cite{AGHMP}. The explicit form of the
height estimate follows Wei \cite[Proposition~5.3]{Wei}, whose zero sum we write out.

\paragraph{The input from Dirichlet $L$-functions.}
We use one consequence of Alp\"oge's theorem \cite{Alpoge}, the Riemann hypothesis for Dirichlet $L$-functions
(formalized in Lean): for every fundamental discriminant $d$,
\[
 L(\sigma,\chi_d)\neq0\qquad\bigl(1-\tfrac1{4\log|d|}\leq\sigma<1\bigr).
 \tag{2.1}
\]
It enters only the height estimate (Proposition~\ref{prop:height}), through Stark's descent of exceptional zeros
\cite[Lemma~9]{Stark}. Since October 2026 there are also unconditional results of this kind. Banks \cite{Banks},
giving a detailed account of an OpenAI manuscript, proves $(1-\beta)\log q\geq3.0\cdot10^{-3}$ for every real zero
$\beta$ of every primitive real $L(s,\chi)$ of conductor $q\geq3$. Cureton and Thorner \cite{CT} prove a zero-free
half-plane $\Re s>31/32$ for every Dedekind zeta function. Liu \cite{Liu} gives an explicit lower bound for $1-\beta$
at every real zero of every Dedekind zeta function. Any of these makes the proof below effective with other
constants (Remark~\ref{rem:other-inputs}). The contribution of this note is therefore not effectivity as such. It is the explicit exponent, the explicit constant, and the
passage to all centre orders.

\paragraph{The neighbouring results, compared.}
\begin{itemize}
\item \emph{Jaswal} \cite[Proposition~3.3]{Jaswal} proves, for the special points of every Shimura variety and under a
zero-free half-plane for the Dedekind zeta function of the splitting field $L$ of the torus, that the image of the
reciprocity map has size at least $c(\delta)\,d_L^{\delta'}$ with $\delta'=1/(16\,t\,s\,n_T(n_T-1))$. This is
Ullmo and Yafaev's bound \cite[Theorem~6.1]{UY}, which assumed GRH, with an explicit constant. Combined with
Cureton--Thorner it is unconditional, and Jaswal uses it to prove an effective Andr\'e--Oort theorem for Shimura
varieties \cite[Theorems~1.1--1.4]{Jaswal}; that theorem is his and we do not claim it. Compared with
Theorem~\ref{thm:main} on $\Ag$, his bound covers every Shimura variety. It is measured in the discriminant of the
Galois closure, and its exponent involves the degree $n_T$ of that closure, which for a CM field of degree $2g$ can
be as large as $2^gg!$. Ours is measured in the discriminant of the centre itself, with exponent $1/(4g)$. The inputs
differ: a zero-free half-plane for the Galois closures there; here a bound for one real zero of quadratic
$L$-functions, an isogeny theorem and the averaged Colmez formula.
\item \emph{Gaudron and R\'emond} \cite[Th\'eor\`eme~1.9(5)]{GR23} bound
$\prod_\ell\disc(\Z_\ell[\mathrm{Im}\,\rho_\ell])$ by $\Xi(A)^2$, where
$\Xi(A)=\bigl((7g)^{8g^2}[K:\Q]\max(1,\log[K:\Q],h_F(A))\bigr)^{2g^2}$. With Lemma~\ref{lem:centre} and
Proposition~\ref{prop:height} below, this already gives an explicit orbit bound, with exponent $1/(4g^2)$. Their
sharper Th\'eor\`eme~1.8 gives the exponent $1/(2g)$ of (1.1). They do not state an orbit bound for CM points, and
their Proposition~13.7, which identifies the algebra generated by the Galois image, is an input to the theorem they
use, not a competing result.
\item \emph{Tsimerman} \cite{Ts12} had proved, earlier, orbit bounds of the form
$(D_K/D_{K^+})^{1/2+o(1)}$ for CM points with Weyl CM field, unconditionally for small $g$ and under GRH for CM fields
in general. \emph{Binyamini, Schmidt and Yafaev} \cite{BSY} reprove Tsimerman's orbit bound for Shimura varieties of
abelian type by point counting and reduce the general case to a height conjecture. \emph{Pila, Shankar and Tsimerman}
\cite{PST} prove the Andr\'e--Oort conjecture in general. These results do not give explicit exponents.
\item For products of modular curves, \emph{Binyamini} \cite{Bin20} proved effective Andr\'e--Oort bounds outside one
possible exceptional imaginary quadratic field.
\end{itemize}

\paragraph{The scope.}
Theorem~\ref{thm:main} is the arithmetic input of the Pila--Zannier strategy for $\Ag$, made explicit. It does not
by itself give an effective Andr\'e--Oort theorem: that also needs an effective count of rational points on the
relevant definable sets, together with the reduction from special subvarieties to special points. For $\Ag$, Jaswal's
Hecke-correspondence argument with Cureton--Thorner's half-plane supplies effective Andr\'e--Oort by another route.
The constants here are far from optimal, and nothing is proved about the true size of the orbits, which for $g=1$ is
$D^{1/2+o(1)}$.

\section{Proof overview}

There are three steps.

\emph{The Galois image lies in the centre.} Let $A/K$ have all its geometric endomorphisms defined over $K$. Each
Galois automorphism acts on the Tate module $T_\ell A$ commuting with every endomorphism. By Faltings's theorem it is
therefore an element of $\End(A)\otimes\Z_\ell$ commuting with all of $\End(A)$, an element of the centre
$R\otimes\Z_\ell$. Moreover the Galois image spans $R\otimes\Q_\ell$. The $\Z_\ell$-algebras $S_\ell$ generated by
the images are thus suborders of $R\otimes\Z_\ell$ of full rank, and the discriminant of $R$ divides the product of
their discriminants (Lemma~\ref{lem:centre}).

\emph{Gaudron--R\'emond bound that product.} Their Th\'eor\`eme~1.5(2) says that the product of the discriminants of
the algebras $S_\ell$ divides an integer $d$ attached to the isogeny class, and Th\'eor\`eme~1.8 bounds $d$ by
$(C_g[K:\Q]\max\{1,\log[K:\Q],h_F(A)+3g/2\})^{2n}$, where $h_F$ is the Faltings height and $n\leq g$ is the
dimension of the product of the distinct simple factors. Everything is now reduced to a height bound of the form
$h_F(A)\ll\log D$.

\emph{The height of a CM abelian variety is at most a multiple of $\log D$.} For CM by a maximal order $\OO_E$, the
averaged Colmez formula expresses the average height over CM types through $-L'(0,\chi)/L(0,\chi)$, where $\chi$ is
the quadratic character of $E/E^+$. The functional equation turns this into a sum over zeros
$\sum\Re\,1/(1-\rho)$. All zeros outside a small square around $1$ contribute $O(\log D_E)$ by a standard
comparison (Lemma~\ref{lem:zero-sum}). A single real zero inside the square would contribute $1/(1-\beta)$, which is
unbounded in terms of $\log D_E$; this is where the ineffectivity of the earlier proofs lies. By Stark's descent such
a zero, if close enough to $1$, is a zero of $L(s,\chi_d)$ for an imaginary quadratic subfield, and (2.1) excludes
it. For a general CM variety we pass by an isogeny of degree at most $[\OO:R]^{2g}$ to one whose endomorphism ring
is maximal (Lemma~\ref{lem:saturation}); the height changes by at most $g\log[\OO:R]\leq\tfrac g2\log D$.

Solving $D\leq(C_gk\,W_g\log D)^{2g}$ for $k=[K:\Q]$ gives (1.2), and the passage from a field of definition of all
endomorphisms to the field of moduli costs the factor $s_g$.

\section{The height of a CM abelian variety}

This section proves the height estimate. For a number field $K$ write $D_K=|\disc K|$. For an abelian variety $A$
over a number field, $h_F(A)$ is Faltings's stable height in his original normalization, and $h_W(A)$ the height
normalized with the archimedean metric $(2\pi)^{-r}\int|\omega\wedge\bar\omega|$; then
$h_W=h_F+\tfrac r2\log\pi$ for $\dim A=r$ \cite[\S2.3]{GR14}.

\begin{lemma}\label{lem:zero-sum}
Let $K\neq\Q$ be a number field, and let $\beta$ be the real zero of $\zeta_K$ in the square
$\{\Re s\geq1-\frac1{4\log D_K},\ |\Im s|\leq\frac1{4\log D_K}\}$, if there is one. Then, summing over the zeros
$\rho$ of $\zeta_K$ in the critical strip with multiplicity,
\[
 0\leq\sum_{\rho\neq\beta}\Re\frac1{1-\rho}<\frac{75}2\log D_K.
\]
\end{lemma}

\begin{proof}
By Stark \cite[Lemmas~3 and~4]{Stark} the square contains at most one zero, and such a zero is real and simple. Put
$t=1/\log D_K$ and $\sigma=1+t$. Every zero $\rho\neq\beta$ lies outside the square, so $|1-\rho|\geq t/4$ and
$|\sigma-\rho|\leq|1-\rho|+t\leq5|1-\rho|$. With $a=1-\Re\rho>0$,
\[
 \frac{\Re(1/(1-\rho))}{\Re(1/(\sigma-\rho))}=\frac a{a+t}\cdot\frac{|\sigma-\rho|^2}{|1-\rho|^2}<25.
\]
The Hadamard product of the completed zeta function gives, for real $1<\sigma<2$,
\[
 \sum_\rho\Re\frac1{\sigma-\rho}=\frac12\log D_K+\frac1\sigma+\frac1{\sigma-1}
 +\frac{\zeta_K'}{\zeta_K}(\sigma)+\frac{r_1}2\Bigl(\psi\bigl(\tfrac\sigma2\bigr)-\log\pi\Bigr)
 +r_2\bigl(\psi(\sigma)-\log2\pi\bigr),
\]
where $(r_1,r_2)$ is the signature and $\psi=\Gamma'/\Gamma$. Here $\zeta_K'/\zeta_K(\sigma)<0$, $\psi(\sigma/2)\leq\psi(1)=-\gamma$ and
$\psi(\sigma)\leq\psi(2)=1-\gamma$. Each real place contributes at most $-\frac12(\gamma+\log\pi)<-0.86$ and each
complex place at most $1-\gamma-\log2\pi<-1.41$. Since $r_1+2r_2\geq2$ and $1/\sigma<1$, the sum of these terms and
$1/\sigma$ is negative. Hence $\sum_\rho\Re(1/(\sigma-\rho))<1/t+\frac12\log D_K=\frac32\log D_K$. All terms are
nonnegative, so the comparison above gives
$\sum_{\rho\neq\beta}\Re(1/(1-\rho))<25\cdot\frac32\log D_K$.
\end{proof}

\begin{proposition}\label{prop:height}
Let $E$ be a CM field of degree $2r$ and $A$ an abelian variety with complex multiplication by $\OO_E$, of any CM
type. Granting (2.1),
\[
 h_F(A)\leq a_r\log D_E+b_r.
\]
\end{proposition}

\begin{proof}
Put $F=E^+$ and $L(s)=\zeta_E(s)/\zeta_F(s)=L(s,\chi_{E/F})$, an entire Hecke $L$-function. Its completion
$\Lambda(s)=(D_E/D_F)^{s/2}\bigl(\pi^{-(s+1)/2}\Gamma(\tfrac{s+1}2)\bigr)^rL(s)$ satisfies $\Lambda(s)=\Lambda(1-s)$,
and its zeros are zeros of $\zeta_E$, with at most the same multiplicity. From the Hadamard product,
$\Re\,\Lambda'/\Lambda(1)=\sum_\rho\Re\,1/(1-\rho)$ over the zeros of $\Lambda$. With the functional equation this gives
\[
 -\frac{L'(0)}{L(0)}=\sum_\rho\Re\frac1{1-\rho}+\frac12\log\frac{D_E}{D_F}+\frac r2\Bigl(\psi\bigl(\tfrac12\bigr)-\log\pi\Bigr).
\]
Apply Lemma~\ref{lem:zero-sum} to $K=E$. All zeros of $L$ other than the exceptional zero $\beta$ of $\zeta_E$
contribute less than $\frac{75}2\log D_E$. If $\beta$ is a zero of $L$ with $\beta\geq1-(16\,r!\log D_E)^{-1}$, then by
Stark's Lemma~9 \cite[pp.~145--146]{Stark} it is a zero of $\zeta_k$ for an imaginary quadratic subfield $k$ of $E$. Since
$\zeta(\sigma)\neq0$ for $0<\sigma<1$, it is a zero of $L(s,\chi_{d_k})$ with $1-\beta\leq1/(16\log D_E)\leq1/(4\log D_k)$.
This is excluded by (2.1). Hence $1/(1-\beta)<16\,r!\log D_E$, and
$-L'(0)/L(0)<(\frac{75}2+16\,r!)\log D_E+\frac12\log(D_E/D_F)+\frac r2(\psi(\frac12)-\log\pi)$.

The averaged Colmez formula \cite[Theorem~1.1]{YZ}, \cite{AGHMP} states
$2^{-r}\sum_\Phi h_W(A_\Phi)=-\frac12L'(0)/L(0)-\frac14\log(D_E/D_F)$, the sum running over the $2^r$ CM types.
Since $\psi(\frac12)<0$, the right side is less than $\frac12(\frac{75}2+16\,r!)\log D_E$. Every abelian variety of
dimension $r$ has $h_W\geq-\frac r2\log2\pi>-4r$ \cite[Corollaire~8.4]{GR14}. Isolating the term of $A$ gives
$h_W(A)<2^{r-1}(\frac{75}2+16\,r!)\log D_E+4r(2^r-1)$, and $h_F\leq h_W$.
\end{proof}

\section{Centre orders and the Galois image}

This section passes from maximal orders to arbitrary ones, and from the height to the Galois image. Let $A$ have
dimension $g$, complex multiplication, and all geometric endomorphisms defined over the number field $K$. Write
$A\sim\prod_{i=1}^tA_i^{m_i}$ with pairwise non-isogenous simple $A_i$, $\End^0(A_i)=E_i$ a CM field of degree
$2r_i$, so $\sum m_ir_i=g$. Put $E=\prod E_i$, $\OO=\prod\OO_{E_i}$, $R=Z(\End A)\subset\OO$, $I=[\OO:R]$ and
$D=|\disc R|=I^2\prod_iD_{E_i}$.

\begin{lemma}\label{lem:saturation}
There is a $K$-isogeny $A\to B$ of degree at most $I^{2g}$ such that $\End B$ is a maximal order. Consequently,
granting (2.1),
\[
 h_F(A)\leq g\bigl(A_g+\tfrac12\bigr)\log D+gB_g.
\]
\end{lemma}

\begin{proof}
Fix $K\subset\C$ and write $A(\C)=V/\Lambda$. Since $I\OO\subset R$, we have $\Lambda\subset\OO\Lambda\subset I^{-1}\Lambda$,
so $[\OO\Lambda:\Lambda]\leq I^{2g}$. The finite subgroup $\OO\Lambda/\Lambda$ is defined over $K$: if $u_j$ is a
$\Z$-basis of $\OO$, it is the sum of the images of $A[I]$ under the $K$-endomorphisms $Iu_j$. Let $B=V/\OO\Lambda$.
Its rational endomorphism algebra is that of $A$, $\prod M_{m_i}(E_i)$, whose elements commute with $E$; hence
$\End B=\End_\OO(\OO\Lambda)$. Each idempotent component of $\OO\Lambda$ is a projective $\OO_{E_i}$-module of rank
$m_i$, whose endomorphism ring is a maximal order of $M_{m_i}(E_i)$. All endomorphisms of $B$ are defined over $K$,
since they are conjugate to those of $A$ by a $K$-isogeny.

R\'emond's theorem \cite[Th\'eor\`eme~1.2]{Remond} decomposes $B$ over $K$ into at most $g$ simple factors with
maximal endomorphism rings, each with CM by some $\OO_{E_i}$. By Proposition~\ref{prop:height}, additivity of the
height, and $D_{E_i}\leq D$, we get $h_F(B)\leq gA_g\log D+gB_g$. Raynaud's isogeny estimate \cite{Raynaud} gives
$|h_F(A)-h_F(B)|\leq\frac12\log\deg(A\to B)\leq g\log I$, and $\log I\leq\frac12\log D$.
\end{proof}

\begin{lemma}\label{lem:centre}
For each prime $\ell$ let $S_\ell=\Z_\ell[\rho_\ell(\Gal(\overline K/K))]\subset\End_{\Z_\ell}(T_\ell A)$. Then
$S_\ell\subset R\otimes\Z_\ell$ and $S_\ell\otimes\Q_\ell=E\otimes\Q_\ell$. Consequently $D$ divides
$\prod_\ell\disc(S_\ell)$, where $\disc(S_\ell)$ is the power of $\ell$ generating the discriminant ideal of the
trace form.
\end{lemma}

\begin{proof}
The commutant of $\End^0(A)\otimes\Q_\ell=\prod M_{m_i}(E_i\otimes\Q_\ell)$ in $\End(V_\ell A)$ is its centre
$E\otimes\Q_\ell$. Galois commutes with all endomorphisms defined over $K$, so its image lies in $E\otimes\Q_\ell$ and
is commutative. Each Galois operator therefore commutes with the whole image and preserves $T_\ell A$. By Faltings's
integral Tate theorem $\End(A)\otimes\Z_\ell=\End_{\Gal}(T_\ell A)$ \cite{Faltings}, each such operator lies in
$\End(A)\otimes\Z_\ell$. Since it commutes with that ring, it lies in its centre, which is $R\otimes\Z_\ell$.
Faltings's semisimplicity and rational Tate theorem, with the double commutant theorem, give
$S_\ell\otimes\Q_\ell=E\otimes\Q_\ell$. So $S_\ell$ has finite index in $R\otimes\Z_\ell$, and
$\disc(S_\ell)=\disc(R\otimes\Z_\ell)\,[R\otimes\Z_\ell:S_\ell]^2$. The $\ell$-part of $D$ is
$\disc(R\otimes\Z_\ell)$.
\end{proof}

\section{Proof of the theorem}

We use three statements of Gaudron and R\'emond \cite{GR23}, applied to the class $\mathcal C$ of abelian varieties
over $K$ isogenous over $K$ to an abelian subvariety of a power of $A$. Let $\mathcal C_0$ be the isogeny class of
$\prod_iA_i$, of dimension $n=\sum_ir_i\leq g$. Let $d$ be the largest discriminant of the endomorphism ring of a member of $\mathcal C_0$ whose
$\ell$-adic Tate modules are free over $\End\otimes\Z_\ell$ \cite[Notation~1.1]{GR23}, and let $\Upsilon$ be the
invariant of \cite[Notation~1.2]{GR23}. Then:
\begin{itemize}
\item $\prod_\ell\disc(S_\ell)$ divides $d$ \cite[Th\'eor\`eme~1.5(2)]{GR23};
\item $d\leq\Upsilon^{2n}$ \cite[after Notation~1.2]{GR23};
\item $\Upsilon\leq241(\mathrm en)^{2n}n^5[K:\Q]\max\{1,\log[K:\Q],h_F(A)+\frac32g\}$
\cite[Th\'eor\`eme~1.8]{GR23}.
\end{itemize}
The discriminants in the first item use the intrinsic trace, as in Lemma~\ref{lem:centre}. The identity behind it is
\cite[Proposition~13.7]{GR23}.

\begin{proof}[Proof of Theorem~\ref{thm:main}]
Put $k=[K:\Q]$ and $M=\max\{1,\log k,h_F(A)+\frac32g\}$. By Lemma~\ref{lem:centre} and the three statements,
$D\leq(C_gkM)^{2g}$, since $n\leq g$ and $C_gkM>1$. Suppose $k\leq D$. Since $D\geq3$, Lemma~\ref{lem:saturation}
gives $h_F(A)+\frac32g\leq g(A_g+\frac12)\log D+(gB_g+\frac32g)\log D/\log3$, so $M\leq W_g\log D$. Hence
$D\leq(C_gkW_g\log D)^{2g}$, which is (1.2). If $k>D$, then (1.2) holds because $C_gW_g\log D>1$ and
$D^{1/(2g)}\leq D$.

Now let $x$ be a CM point. Adjoin to the principally polarized variety $A_x$ a basis of its $3$-torsion, and a
primitive cube root of unity. A polarization-preserving automorphism fixing the $3$-torsion is the identity, because
the principal congruence subgroup of level three is torsion-free, so the marked object descends to its field of
moduli. This field has degree at most $s_g$ over $\Q(x)$, since there are fewer than $3^{4g^2}$ bases. Every
geometric endomorphism $u$ is defined over it: for $\sigma$ in its Galois group, $\sigma u-u$ vanishes on $A[3]$, so
$\sigma u-u\in3\End(A)$. The finite image of the Galois group in $\mathrm{GL}(\End A)$ therefore lies in the
level-three congruence subgroup, and is trivial. Applying (1.2) over this field gives the first inequality in (1.1).
For the second, use $\log D\leq\frac{4g}{\mathrm e}D^{1/(4g)}$, the maximum of $te^{-t/(4g)}$ being $4g/\mathrm e$.
\end{proof}

\begin{proof}[Proof of Corollary~\ref{cor:surfaces}]
For $g=2$, $s_2=2\cdot3^{16}<10^8$, $C_2=241\cdot(2\mathrm e)^4\cdot32<10^7$, and $W_2=279+52/\log3<331$. For a
rational point, (1.1) gives $D^{1/4}/\log D<331\cdot10^{15}$. The function $e^{t/4}/t$ increases for $t>4$, and at
$t=\log10^{80}<200$ it exceeds $10^{20}/200$, which is larger. So $D<10^{80}$.
\end{proof}

\begin{remark}\label{rem:other-inputs}
The hypothesis (2.1) is used only in the proof of Proposition~\ref{prop:height}, to bound $1/(1-\beta)$ for a real zero
$\beta$ of $L(s,\chi_{d_k})$ with $k$ an imaginary quadratic subfield of $E$. Suppose instead that every real zero of
every such $L$-function satisfies $(1-\beta)\log D_k\geq\delta$. Then $1-\beta\geq\delta/\log D_E$, and the proof gives
Theorem~\ref{thm:main} with $a_r$ replaced by $2^{r-1}(\frac{75}2+\max\{16\,r!,1/\delta\})$. Banks's bound
\cite[abstract]{Banks} has $\delta=3.0\cdot10^{-3}$, so $1/\delta<334$. This changes $a_1,a_2,a_3$ and
leaves $a_r$ unchanged for $r\geq4$. Under (2.1), $\delta=\frac14$ suffices, and $16\,r!\geq16>4$.
\end{remark}

\section{What remains open}

The exponent $1/(4g)$ is an artefact of the method: the isogeny theorem loses the power $2n$, and the heights are
bounded crudely. For $g=1$ the truth is $D^{1/2+o(1)}$, and for CM points with Weyl CM field Tsimerman's conditional
bound \cite{Ts12} has exponent $1/2$ in $D_K/D_{K^+}$. An explicit exponent independent of $g$, or even $1/(2g)$
without the logarithm, would require a different input. Making the whole Andr\'e--Oort theorem for $\Ag$ explicit
along the Pila--Zannier route needs, besides Theorem~\ref{thm:main}, an effective point count for the uniformization
of $\Ag$ on a fundamental domain. Binyamini's work on Log-Noetherian functions \cite[\S10]{BinLN} identifies the
remaining analytic data.

\section*{Method}

The theorem and its proofs were produced by GPT-6 Astra (ChatGPT Pro mode). Each statement was re-proved from its
statement by separate workers of the same model. Claude Opus 5.5 at effort high checked every proof step by step and
recomputed the constants and the discriminant cutoff. Claude Opus 5.5 at effort max read Jaswal's paper and Gaudron--R\'emond's at the
places compared here. This text was written by Claude Opus 5.5 at effort high. It read at the source
Gaudron--R\'emond's Notations~1.1--1.2, Th\'eor\`emes~1.5, 1.8, 1.9 and Proposition~13.7 (authors' preprint of
September 2022) and Tsimerman's \cite{Ts18} (arXiv version~5). For Banks, Liu and Binyamini--Schmidt--Yafaev it read
the abstracts. The other citations follow the locators of the producing model's checked sources. No human has
checked this work. Our text is dedicated to the public domain under CC0~1.0. Records: \url{\workbenchurl}; permanent
record: \url{\recordurl}.

\begin{thebibliography}{99}
\small\raggedright
\bibitem{Alpoge} L.~Alp\"oge, \emph{Remark}, 20 November 2026, \url{https://alpo.ge/remark.pdf}; Lean archive
\url{https://github.com/CocoSgt/riemann-hypothesis-lean}.
\bibitem{AGHMP} F.~Andreatta, E.~Z.~Goren, B.~Howard, K.~Madapusi Pera, Faltings heights of abelian varieties with
complex multiplication, Ann.\ of Math.\ 187 (2018), 391--531.
\bibitem{Banks} W.~D.~Banks, On the OpenAI proof that Landau--Siegel zeros do not exist, arXiv:2610.15483 (2026).
\bibitem{Bin20} G.~Binyamini, Some effective estimates for Andr\'e--Oort in $Y(1)^n$, with an appendix by
E.~Kowalski, J.\ reine angew.\ Math.\ 767 (2020), 17--35.
\bibitem{BinLN} G.~Binyamini, Log-Noetherian functions, arXiv:2405.16963.
\bibitem{BSY} G.~Binyamini, H.~Schmidt, A.~Yafaev, Lower bounds for Galois orbits of special points on Shimura
varieties: a point-counting approach, Math.\ Ann.\ (2023); arXiv:2104.05842.
\bibitem{CT} D.~Cureton, J.~Thorner, A quasi-Riemann hypothesis for Dedekind zeta functions, hexagon:2610.00200v2
(2026), Theorem~1.1.
\bibitem{Faltings} G.~Faltings, Endlichkeitss\"atze f\"ur abelsche Variet\"aten \"uber Zahlk\"orpern, Invent.\ Math.\
73 (1983), 349--366.
\bibitem{GR14} \'E.~Gaudron, G.~R\'emond, Th\'eor\`eme des p\'eriodes et degr\'es minimaux d'isog\'enies, Comment.\
Math.\ Helv.\ 89 (2014), 343--403.
\bibitem{GR23} \'E.~Gaudron, G.~R\'emond, Nouveaux th\'eor\`emes d'isog\'enie, M\'em.\ Soc.\ Math.\ France 176 (2023);
authors' preprint (September 2022), \url{https://lmbp.uca.fr/~gaudron/art20.pdf}, pp.~10--15 and \S13.
\bibitem{Jaswal} M.~Jaswal, An effective Andr\'e--Oort theorem for Shimura varieties, hexagon:2611.00339v3 (2026).
\bibitem{Liu} B.~Liu, Landau--Siegel zeros over varying number fields, hexagon:2611.00080v2 (2026).
\bibitem{PST} J.~Pila, A.~N.~Shankar, J.~Tsimerman, with an appendix by H.~Esnault and M.~Groechenig, Canonical
heights on Shimura varieties and the Andr\'e--Oort conjecture, arXiv:2109.08788.
\bibitem{Raynaud} M.~Raynaud, Hauteurs et isog\'enies, Ast\'erisque 183 (1990), 199--234.
\bibitem{Remond} G.~R\'emond, Vari\'et\'es ab\'eliennes et ordres maximaux, Rev.\ Mat.\ Iberoam.\ 33 (2017),
1173--1195.
\bibitem{Stark} H.~M.~Stark, Some effective cases of the Brauer--Siegel theorem, Invent.\ Math.\ 23 (1974), 135--152.
\bibitem{Ts12} J.~Tsimerman, Brauer--Siegel for arithmetic tori and lower bounds for Galois orbits of special points,
J.\ Amer.\ Math.\ Soc.\ 25 (2012), 1091--1117; arXiv:1103.5619.
\bibitem{Ts18} J.~Tsimerman, The Andr\'e--Oort conjecture for $\mathcal A_g$, Ann.\ of Math.\ 187 (2018), 379--390;
arXiv:1506.01466v5, Theorems~1.1, 4.1, 5.1, 5.2.
\bibitem{UY} E.~Ullmo, A.~Yafaev, Nombre de classes des tores de multiplication complexe et bornes inf\'erieures pour
les orbites galoisiennes de points sp\'eciaux, Bull.\ Soc.\ Math.\ France 143 (2015), 197--228.
\bibitem{Wei} X.~Wei, Smallness of Faltings heights of CM abelian varieties, J.\ Number Theory 242 (2023); arXiv:2111.00617,
Proposition~5.3.
\bibitem{YZ} X.~Yuan, S.-W.~Zhang, On the averaged Colmez conjecture, Ann.\ of Math.\ 187 (2018), 533--638.
\end{thebibliography}
\end{document}
